% =====================================================================
%  setup/macros.tex -- notation and reference helpers.
%  One way to write each thing; see frontmatter/how_to_read.tex (Notation).
% =====================================================================

% ---------------- Number sets ----------------
\newcommand{\N}{\mathbb{N}}
\newcommand{\Z}{\mathbb{Z}}
\newcommand{\R}{\mathbb{R}}

% ---------------- Delimiters (starred forms auto-size) ----------------
\DeclarePairedDelimiter{\floor}{\lfloor}{\rfloor}
\DeclarePairedDelimiter{\ceil}{\lceil}{\rceil}
\DeclarePairedDelimiter{\abs}{\lvert}{\rvert}
\DeclarePairedDelimiter{\set}{\{}{\}}
\DeclarePairedDelimiter{\card}{\lvert}{\rvert}          % cardinality |S|

% ---------------- Tree operators (roman, like \gcd) ----------------
% Each takes a node or a tree in parentheses: \depth(v), \height(T), ...
\DeclareMathOperator{\depth}{depth}
\DeclareMathOperator{\height}{height}
\DeclareMathOperator{\size}{size}                        % subtree node count
\DeclareMathOperator{\level}{level}
\DeclareMathOperator{\degop}{deg}                        % degree of a node
\DeclareMathOperator{\parent}{parent}
\DeclareMathOperator{\children}{children}
\DeclareMathOperator{\leaves}{leaves}
\DeclareMathOperator{\subtree}{subtree}
\DeclareMathOperator{\ancestors}{ancestors}
\DeclareMathOperator{\lca}{lca}                          % lowest common ancestor
\DeclareMathOperator{\diam}{diam}                        % diameter (longest path)
\DeclareMathOperator{\dist}{dist}                        % distance between nodes
\DeclareMathOperator{\bal}{bal}                          % balance factor
\DeclareMathOperator{\tin}{tin}                          % Euler-tour entry time
\DeclareMathOperator{\tout}{tout}                        % Euler-tour exit time

% ---------------- Text helpers ----------------
\newcommand{\code}[1]{\texttt{#1}}                            % inline code; escape _ & % # as usual
\newcommand{\cpp}{C\nolinebreak\hspace{-.05em}\raisebox{.25ex}{\scriptsize\bfseries ++}}
\newcommand{\term}[1]{\textbf{\textcolor{deepblue}{#1}}\index{#1}}   % first use of a key term
\newcommand{\nil}{\code{null}}                                % the empty child pointer

% ---------------- Complexity ----------------
\newcommand{\bigO}[1]{\ensuremath{O(#1)}}
\newcommand{\bigTheta}[1]{\ensuremath{\Theta(#1)}}

% ---------------- Cross-references (label prefixes are fixed) ----------------
%  ch:   chapter          sec:  topic (section)     def:/thm:/lem:/cor:/prop:/ex:  boxes
%  fig:  figure           tab:  table               lst:  code listing
%  prob: solved problem   eq:   equation
%  Every label after the prefix starts with the chapter tag, e.g. thm:c01-edges-nodes.
\newcommand{\chref}[1]{Chapter~\ref{ch:#1}}
\newcommand{\secref}[1]{\S\ref{sec:#1}}
\newcommand{\defref}[1]{Definition~\ref{def:#1}}
\newcommand{\thmref}[1]{Theorem~\ref{thm:#1}}
\newcommand{\lemref}[1]{Lemma~\ref{lem:#1}}
\newcommand{\corref}[1]{Corollary~\ref{cor:#1}}
\newcommand{\propref}[1]{Proposition~\ref{prop:#1}}
\newcommand{\exref}[1]{Example~\ref{ex:#1}}
\newcommand{\figref}[1]{Figure~\ref{fig:#1}}
\newcommand{\tabref}[1]{Table~\ref{tab:#1}}
\newcommand{\coderef}[1]{Code~\ref{lst:#1}}
\newcommand{\probref}[1]{W\ref{prob:#1}}
